\chapter{Обзор используемых алгоритмов оптимизации} \label{ch2}
	
% не рекомендуется использовать отдельную section <<введение>> после лета 2020 года
%\section{Введение} \label{ch2:intro}

В данной главе приведена краткая справка по используемым в работе алгоритмам решения поставленной задачи \labelcref{eq:goal-function}
\section{Генетический алгоритм} \label{ch2:title-abbr} %название по-русски
Генетические алгоритмы (ГА) - класс методов оптимизации, вдохновленных принципами естественного отбора и генетики.\\
Основная идея ГА заключается в том, чтобы эмулировать естественный процесс эволюции, в котором лучшие особи имеют больше шансов передать свои характеристики следующему поколению.
Путем применения операторов скрещивания, мутации и отбора, ГА итеративно улучшают текущее поколение решений, стремясь к нахождению оптимального или приближенного к нему решения.
\newpage
\subsection{Псевдокод генетического алгоритма}
		\begin{algorithm} %[h]
		\SetKwFunction{algoDTestsFDSCALING}{} 
		\SetKwProg{myalg}{Algorithm}{}{} %write in 2nd agrument <<Algorithm>>, <<Procedure>> etc
		\nonl\myalg{\algoDTestsFDSCALING}{
			\KwInput{\\$n$ - количество поколений,\\
					   $p$ - количество скрещивающихся хромосом в поколении,\\
					   $q$ - количество хромосом в поколении,\\
			 		   $[t_{min}, t_{max}]$ диапазон возможных значений времени фазы.
			} 
			\KwOutput{Оптимальное значение $f(X_{*})$.} 
			Родители $\leftarrow$ $\{\text{Сгенерированная случайным образом популяция}\}$\;
			\While {не выполнено условие остановки} {
				Вычислить значение целевой функции для каждой хромосомы-родителя в популяции\;
				Дети $\leftarrow \emptyset$\;
				\While {|Дети| < |Родители|} {
					На основании значения ЦФ выбрать пару родителей для скрещивания\;
					Скрестить родителей, чтобы создать детей $c_1$ и $c_2$\;
					Дети $\leftarrow$ Дети $\cup \ \{c_1, c_2\}$\;
				}
				Вызвать сл. образом мутацию детей\;
				Родители $\leftarrow$ Дети\;
			}		
		}
		\caption{Псевдокод генетического алгоритма}
	\end{algorithm} 
	 % пример оформления псевдокода алгоритма 	
\section{Метод роя частиц}
\subsection{Псевдокод метода роя частиц}
\section{Выводы} \label{ch2:conclusion}


%% Вспомогательные команды - Additional commands
%
%\newpage % принудительное начало с новой страницы, использовать только в конце раздела
%\clearpage % осуществляется пакетом <<placeins>> в пределах секций
%\newpage\leavevmode\thispagestyle{empty}\newpage % 100 % начало новой страницы